\section{Background}

Machine-learning-based classification and clustering are now routine in biological and biomedical research, yet the profiling experiments that feed these analyses now routinely contain many more measured features than samples. The resulting ``curse of dimensionality'' can inflate apparent performance, drive overfitting, and destabilize model selection, so practitioners combine dimensionality reduction, cross-validation, regularization, and permutation-based significance estimation in order to obtain trustworthy predictions \cite{harris2020numpy,virtanen2020scipy,pedregosa2011scikit}. Permutation testing---in which the target label is repeatedly scrambled across samples and the model is retrained---provides a direct, nonparametric way to characterize the distribution of performance expected by chance, against which the model trained on actual labels can be compared \cite{ojala2010permutation,breiman1996bagging}.

An important but comparatively less well-served subfield of semi-supervised classification is Positive-Unlabeled (PU) learning. PU learning targets the common biomedical regime in which a small number of confirmed positive instances are known but no confirmed negatives are available; the remaining samples are pooled into an ``unlabeled'' set that is presumed to contain both positives and negatives \cite{bekker2020pulearning,elkan2008learning,li2022pulearning}. PU learning has been successfully applied to disease-gene identification \cite{yang2014epu}, miRNA-disease association prediction \cite{wang2019nsemda}, and protease substrate prediction \cite{li2019protease}. Two broad design patterns dominate: Reliable Negative Selection, which identifies a confident subset of true negatives from the unlabeled pool, and Base Classifier Adaptation, which directly modifies the learning procedure to accommodate unlabeled data \cite{li2022pulearning}.

Among Base Classifier Adaptation methods, the transductive PU Bagging ensemble of Mordelet and Vert constructs many base classifiers by sampling the unlabeled set with replacement and temporarily labeling the bootstrapped subset as negative, then aggregates the out-of-bag class-1 probabilities across replicates \cite{mordelet2014bagging,breiman1996bagging}. PU Bagging is relatively insensitive to the choice of base classifier and, in high-dimensional biological data, can outperform a single biased support vector machine. Within Reliable Negative Selection, Liu et al.\ introduced the ``spy'' technique, in which a small random subset of the known positives is mixed into the unlabeled set and used to calibrate a class-1 probability threshold that selects reliable negatives \cite{liu2003spy}. Bekker and Davis later argued that the spy procedure is less reliable when the known-positive set is small, but spies remain a widely used calibration strategy \cite{bekker2020pulearning}.

The fundamental difficulty in deploying any PU learning method is that validation metrics such as accuracy, F1 score, Matthews correlation coefficient, and area under the receiver operating characteristic curve all presuppose a ground-truth negative class \cite{hanley1982auc}. In many biomedical use cases, definitive negative labels are either prohibitively costly---requiring exhaustive gold-standard assays---or biologically undefined, as when a gene is only ``currently not known to be associated with a disease'' rather than confidently disease-unrelated \cite{yang2014epu,wang2019nsemda}. A limited body of work has attempted to validate PU predictions using positive examples alone: in particular, Explicit Positive Recall (EPR) scores a model by the fraction of known positives that receive a predicted label of positive \cite{cheng2015epr}. Whether high recall of known positives acts as a reliable surrogate for reliable prediction of the unlabeled set is, however, context-dependent and, in general, ungrounded without a null reference.

Here, we propose a generalizable methodology to evaluate the robustness of positive-set scores in transductive PU learning without relying on negative examples. We integrate a k-fold spy-positive design \cite{liu2003spy} with the PU Bagging classifier of Mordelet and Vert \cite{mordelet2014bagging}, using a support vector machine with a radial basis function kernel \cite{pedregosa2011scikit} as the base learner, and we score each positive sample using both EPR \cite{cheng2015epr} and a Mean Bagging Score (MBS). We then compare the actual known-positive score distribution against the distribution obtained after repeated label permutations, using Welch's t-test, Cliff's delta effect size with a 95\% confidence interval \cite{cliff1993dominance}, and a one-sample z-test. We validate this pipeline on synthetic high-dimensional datasets generated with scikit-learn \cite{pedregosa2011scikit,abraham2014scikit}, on the Wisconsin Diagnostic Breast Cancer (WDBC) dataset \cite{wolberg1999nuclear,wolberg2002computer}, on a BRCA vs.\ LUAD subset of The Cancer Genome Atlas PANCAN cohort \cite{weinstein2013tcga}, and on a humoral immune-response dataset from a rhesus macaque HIV vaccine study \cite{lakhashe2022virosomal,marasini2021enshivig,malherbe2021adenovirus}. The remainder of the manuscript describes the pipeline, reports results across varied class separations, class-ratio imbalances, and known-positive proportions, and documents the stability of the comparison metrics as a function of the number of permutation replicates.
